% GENERATED FILE. Edit canonical lessons, then run npm run generate:books.
% canonical-source-hash: fnv1a64:ca839e77a9e34e9e
% canonical-lessons: TE-C229-tapakayalu, TE-C229-gangireddu, TE-C229-enda, TE-C229-cinuku, TE-C229-pogamancu

\chapter{Firecrackers, Decorated performing bull seen at Sankranti, Sunshine, Raindrop, Fog}
\label{ch:te-firecrackers-decorated-performing-bull-seen-at-sankranti-sunshine-raindrop-fog}
\hlchaptermodality{229}

\begin{chapteropening}
\textbf{By the end of this chapter:} \emph{I can say firecrackers, a decorated performing bull seen at Sankranti, sunshine, a raindrop and fog.}

Complete the last lesson of chapter 229: I can say firecrackers, a decorated performing bull seen at Sankranti, sunshine, a raindrop and fog.
\end{chapteropening}

\section[\texorpdfstring{\te{టపాకాయలు} (ṭapākāyalu)}{టపాకాయలు (ṭapākāyalu)}]{\te{టపాకాయలు} (ṭapākāyalu) — firecrackers}

\label{lesson:TE-C229-tapakayalu}



\begin{quote}
Before the new one: say the Telugu for a nest, then the Telugu for a tail.
\end{quote}

\subsection*{You'll want to know: \te{టపాకాయలు}}
\textbf{\te{టపాకాయలు}} — \emph{ṭapākāyalu} — \textquotedblleft{}firecrackers\textquotedblright{}.

\begin{grammarlens}[title={What you've built}]
One more word for this chapter.
\end{grammarlens}

\subsection*{Your turn}
\begin{itemize}
\raggedright
  \item \emph{Say it:} \emph{ṭapākāyalu}
  \item \emph{Say it:} \emph{ṭapākāyalu}, once more
  \item \emph{Say it:} say \emph{ṭapākāyalu}
  \item \emph{Recall:} say the Telugu for a nest, then the Telugu for a tail, then say \emph{ṭapākāyalu} again
\end{itemize}

\begin{tcolorbox}[breakable,colback=teal!4,colframe=teal!35!black,title={Before you move on}]
What does \te{టపాకాయలు} mean? (\textquotedblleft{}Firecrackers\textquotedblright{}.) Say it once more.
\end{tcolorbox}
\section[\texorpdfstring{\te{గంగిరెద్దు} (gaṅgireddu)}{గంగిరెద్దు (gaṅgireddu)}]{\te{గంగిరెద్దు} (gaṅgireddu) — a decorated performing bull seen at Sankranti}

\label{lesson:TE-C229-gangireddu}



\begin{quote}
Before the new one: say the Telugu for a tail, then the Telugu for firecrackers.
\end{quote}

\subsection*{You'll want to know: \te{గంగిరెద్దు}}
\textbf{\te{గంగిరెద్దు}} — \emph{gaṅgireddu} — \textquotedblleft{}a decorated performing bull seen at Sankranti\textquotedblright{}.

\begin{grammarlens}[title={What you've built}]
One more word for this chapter.
\end{grammarlens}

\subsection*{Your turn}
\begin{itemize}
\raggedright
  \item \emph{Say it:} \emph{gaṅgireddu}
  \item \emph{Say it:} \emph{gaṅgireddu}, once more
  \item \emph{Say it:} say \emph{gaṅgireddu}
  \item \emph{Recall:} say the Telugu for a tail, then the Telugu for firecrackers, then say \emph{gaṅgireddu} again
\end{itemize}

\begin{tcolorbox}[breakable,colback=teal!4,colframe=teal!35!black,title={Before you move on}]
What does \te{గంగిరెద్దు} mean? (\textquotedblleft{}A decorated performing bull seen at Sankranti\textquotedblright{}.) Say it once more.
\end{tcolorbox}
\section[\texorpdfstring{\te{ఎండ} (eṇḍa)}{ఎండ (eṇḍa)}]{\te{ఎండ} (eṇḍa) — sunshine, the sun's heat}

\label{lesson:TE-C229-enda}



\begin{quote}
Before the new one: say the Telugu for firecrackers, then the Telugu for a decorated performing bull seen at Sankranti.
\end{quote}

\subsection*{You'll want to know: \te{ఎండ}}
\textbf{\te{ఎండ}} — \emph{eṇḍa} — \textquotedblleft{}sunshine, the sun's heat\textquotedblright{}.

\begin{grammarlens}[title={What you've built}]
One more word for this chapter.
\end{grammarlens}

\subsection*{Your turn}
\begin{itemize}
\raggedright
  \item \emph{Say it:} \emph{eṇḍa}
  \item \emph{Say it:} \emph{eṇḍa}, once more
  \item \emph{Say it:} say \emph{eṇḍa}
  \item \emph{Recall:} say the Telugu for firecrackers, then the Telugu for a decorated performing bull seen at Sankranti, then say \emph{eṇḍa} again
\end{itemize}

\begin{tcolorbox}[breakable,colback=teal!4,colframe=teal!35!black,title={Before you move on}]
What does \te{ఎండ} mean? (\textquotedblleft{}Sunshine, the sun's heat\textquotedblright{}.) Say it once more.
\end{tcolorbox}
\section[\texorpdfstring{\te{చినుకు} (cinuku)}{చినుకు (cinuku)}]{\te{చినుకు} (cinuku) — a raindrop}

\label{lesson:TE-C229-cinuku}



\begin{quote}
Before the new one: say the Telugu for a decorated performing bull seen at Sankranti, then the Telugu for sunshine.
\end{quote}

\subsection*{You'll want to know: \te{చినుకు}}
\textbf{\te{చినుకు}} — \emph{cinuku} — \textquotedblleft{}a raindrop\textquotedblright{}.

\begin{grammarlens}[title={What you've built}]
One more word for this chapter.
\end{grammarlens}

\subsection*{Your turn}
\begin{itemize}
\raggedright
  \item \emph{Say it:} \emph{cinuku}
  \item \emph{Say it:} \emph{cinuku}, once more
  \item \emph{Say it:} say \emph{cinuku}
  \item \emph{Recall:} say the Telugu for a decorated performing bull seen at Sankranti, then the Telugu for sunshine, then say \emph{cinuku} again
\end{itemize}

\begin{tcolorbox}[breakable,colback=teal!4,colframe=teal!35!black,title={Before you move on}]
What does \te{చినుకు} mean? (\textquotedblleft{}A raindrop\textquotedblright{}.) Say it once more.
\end{tcolorbox}
\section[\texorpdfstring{\te{పొగమంచు} (pogamañcu)}{పొగమంచు (pogamañcu)}]{\te{పొగమంచు} (pogamañcu) — fog}

\label{lesson:TE-C229-pogamancu}



\begin{quote}
Before the new one: say the Telugu for sunshine, then the Telugu for a raindrop.
\end{quote}

\subsection*{You'll want to know: \te{పొగమంచు}}
\textbf{\te{పొగమంచు}} — \emph{pogamañcu} — \textquotedblleft{}fog\textquotedblright{}.

\begin{grammarlens}[title={What you've built}]
One more word for this chapter.
\end{grammarlens}

\subsection*{Your turn}
\begin{itemize}
\raggedright
  \item \emph{Say it:} \emph{pogamañcu}
  \item \emph{Say it:} \emph{pogamañcu}, once more
  \item \emph{Say it:} say \emph{pogamañcu}
  \item \emph{Recall:} say the Telugu for sunshine, then the Telugu for a raindrop, then say \emph{pogamañcu} again
\end{itemize}

\begin{tcolorbox}[breakable,colback=teal!4,colframe=teal!35!black,title={Before you move on}]
What does \te{పొగమంచు} mean? (\textquotedblleft{}Fog\textquotedblright{}.) Say it once more.
\end{tcolorbox}
